% Standalone problem document, generated from the adjacent README.md.
% Compile with XeLaTeX. Mathematical content is maintained in Markdown.
\documentclass[11pt,a4paper]{article}
\usepackage[fontsize=10.5pt]{scrextend}
\usepackage[margin=22mm,headheight=15pt,headsep=9mm,footskip=11mm]{geometry}
\usepackage{fontspec}
\setmainfont{texgyrepagella}[Extension=.otf,UprightFont=*-regular,BoldFont=*-bold,ItalicFont=*-italic,BoldItalicFont=*-bolditalic]
\setsansfont{texgyreheros}[Extension=.otf,UprightFont=*-regular,BoldFont=*-bold,ItalicFont=*-italic,BoldItalicFont=*-bolditalic]
\setmonofont{texgyrecursor}[Extension=.otf,UprightFont=*-regular,BoldFont=*-bold,ItalicFont=*-italic,BoldItalicFont=*-bolditalic,Scale=MatchLowercase]
\usepackage{amsmath,amssymb}
\usepackage{unicode-math}
\setmathfont{texgyrepagella-math.otf}
\usepackage{xcolor,microtype,enumitem,needspace,fancyhdr}
\usepackage{bookmark}
\definecolor{ink}{HTML}{17344B}
\definecolor{accent}{HTML}{197D80}
\definecolor{muted}{HTML}{596773}
\hypersetup{colorlinks=true,linkcolor=accent,urlcolor=accent,pdftitle={IV-05 - Exact
solution hulls for inverse M-matrix
intervals},pdfauthor={Open Problems in Numerical Linear Algebra}}
\urlstyle{same}
\setlength{\parindent}{0pt}
\setlength{\parskip}{5pt plus 1pt minus 1pt}
\linespread{1.04}
\setlength{\emergencystretch}{3em}
\widowpenalty=10000
\clubpenalty=10000
\displaywidowpenalty=10000
\allowdisplaybreaks[1]
\setlist{nosep,leftmargin=1.5em}
\providecommand{\tightlist}{\setlength{\itemsep}{0pt}\setlength{\parskip}{0pt}}
\providecommand{\pandocbounded}[1]{#1}
\pagestyle{fancy}
\fancyhf{}
\fancyhead[L]{\sffamily\footnotesize\color{muted}OPEN PROBLEMS IN NUMERICAL LINEAR ALGEBRA}
\fancyhead[R]{\sffamily\bfseries\footnotesize\color{accent}IV-05}
\fancyfoot[L]{\parbox[b]{0.9\textwidth}{\sffamily\fontsize{8}{10}\selectfont\color{muted}Editorial ratings; a bounded search cannot certify that a problem remains unsolved.\\Literature check: 2026-09-10}}
\fancyfoot[R]{\sffamily\footnotesize\color{muted}\thepage}
\renewcommand{\headrulewidth}{0.3pt}
\renewcommand{\headrule}{\hbox to\headwidth{\color{accent}\leaders\hrule height \headrulewidth\hfill}}
\makeatletter
\renewcommand{\section}{\Needspace{5\baselineskip}\@startsection{section}{1}{0pt}{12pt plus 2pt minus 2pt}{4pt}{\sffamily\bfseries\large\color{ink}}}
\renewcommand{\subsection}{\Needspace{4\baselineskip}\@startsection{subsection}{2}{0pt}{9pt plus 1pt minus 1pt}{3pt}{\sffamily\bfseries\normalsize\color{ink}}}
\makeatother
\setcounter{secnumdepth}{0}
\begin{document}
{\sffamily\small\color{muted}Intervals and absolute value equations\par}
\vspace{3pt}
{\raggedright\hyphenpenalty=10000\exhyphenpenalty=10000\sffamily\bfseries\fontsize{21}{24}\selectfont\color{ink}Exact
solution hulls for inverse M-matrix intervals\par}
\vspace{7pt}
{\sffamily\small\textbf{Difficulty:} challenging\quad\textbf{Importance:} interesting
to specialist\par}
{\sffamily\small\bfseries\color{accent}Status: Open\par}
\vspace{4pt}
{\color{accent}\hrule height 0.8pt}
\vspace{6pt}
\textbf{Rating rationale:} A polynomial hull algorithm would need to
exploit the inverse-M promise beyond the known endpoint reductions,
warranting challenging. Its direct impact is on a specialist class of
verified linear-system computations.

\textbf{Area:} verified linear systems; interval complexity

\subsection{Problem statement}\label{problem-statement}

For arbitrary \(n\ge1\), the input consists of rational matrices
\(L\le U\) and rational vectors \(l\le u\), with entrywise inequalities.
Let

\[
\mathcal A=\{A\in\mathbb R^{n\times n}:L\le A\le U\},\qquad
\mathcal b=\{b\in\mathbb R^n:l\le b\le u\}.
\]

Assume as a promise that every \(A\in\mathcal A\) is an inverse
M-matrix: \(A^{-1}\) has nonpositive off-diagonal entries and is a
nonsingular M-matrix. Here a nonsingular M-matrix means an invertible
real matrix with nonpositive off-diagonal entries and entrywise
nonnegative inverse. The entries of \(A\) and \(b\) vary independently.

Is there a deterministic algorithm that returns the exact rational
endpoints of

\[
\prod_{i=1}^n
\left[\min_{A\in\mathcal A,\ b\in\mathcal b}(A^{-1}b)_i,
      \max_{A\in\mathcal A,\ b\in\mathcal b}(A^{-1}b)_i\right]
\]

in time polynomial in the binary input length? Behavior outside the
promise is unrestricted, and checking the promise is not part of the
requested algorithm. This is the precise bit-complexity formulation of
the source's efficient interval-system solution-hull question.

The promise makes the solution set nonempty and compact. The aim is its
smallest coordinatewise enclosure. Computing the interval hull of
\(A^{-1}\) separately is already easier in this matrix class;
multiplying that enclosure by \(\mathcal b\) need not capture the
dependence among entries of the inverse.

\newpage
\subsection{References}\label{references}

Milan Hladík,
\href{https://doi.org/10.1007/978-3-030-31041-7_16}{\emph{An overview of
polynomially computable characteristics of special interval matrices}},
in \emph{Beyond Traditional Probabilistic Data Processing Techniques},
Springer (2020), 295--310;
\href{https://arxiv.org/pdf/1711.08732}{arXiv:1711.08732v1}, §9, p.~11,
question immediately before Theorem 25. Theorem 24 concerns the
different inverse-matrix hull; Theorem 25 specifies the right-hand-side
vertices attaining each endpoint.

\subsection{Status check}\label{status-check}

On 2026-09-10, searches for inverse M-matrix interval systems, exact
solution hull, polynomial/efficient computation, and 2025/2026 found no
resolution. The author's publication list through 2026 and
Garloff--Al-Saafin--Adm's
\href{https://reliable-computing.org/reliable-computing-28-pp-056-070.pdf}{2021
interval-property paper}, §4, were checked. The latter addresses
recognizing the matrix class, not this promised hull computation. The
explicit openness statement remains historical.

\subsection{Independent audit ---
2026-09-10}\label{independent-audit-2026-09-10}

The full Hladík preprint, §9, asks the efficient hull question
immediately before Theorem 25. That theorem fixes the optimizing
right-hand-side vertices but leaves the matrix optimization unresolved;
it is not a polynomial algorithm or a solved subclass of this target.
The 2021 interval-property paper addresses recognition instead. Later
searches specific to inverse M-matrix solution hulls found no
resolution; the explicit open-status evidence remains historical.
\end{document}
